% Options for packages loaded elsewhere
\PassOptionsToPackage{unicode}{hyperref}
\PassOptionsToPackage{hyphens}{url}
\PassOptionsToPackage{dvipsnames,svgnames,x11names}{xcolor}
\documentclass[
  11pt,
]{article}
\usepackage{xcolor}
\usepackage[margin=2.4cm]{geometry}
\usepackage{amsmath,amssymb}
\setcounter{secnumdepth}{-\maxdimen} % remove section numbering
\usepackage{iftex}
\ifPDFTeX
  \usepackage[T1]{fontenc}
  \usepackage[utf8]{inputenc}
  \usepackage{textcomp} % provide euro and other symbols
\else % if luatex or xetex
  \usepackage{unicode-math} % this also loads fontspec
  \defaultfontfeatures{Scale=MatchLowercase}
  \defaultfontfeatures[\rmfamily]{Ligatures=TeX,Scale=1}
\fi
\usepackage{lmodern}
\ifPDFTeX\else
  % xetex/luatex font selection
\fi
% Use upquote if available, for straight quotes in verbatim environments
\IfFileExists{upquote.sty}{\usepackage{upquote}}{}
\IfFileExists{microtype.sty}{% use microtype if available
  \usepackage[]{microtype}
  \UseMicrotypeSet[protrusion]{basicmath} % disable protrusion for tt fonts
}{}
\usepackage{setspace}
\makeatletter
\@ifundefined{KOMAClassName}{% if non-KOMA class
  \IfFileExists{parskip.sty}{%
    \usepackage{parskip}
  }{% else
    \setlength{\parindent}{0pt}
    \setlength{\parskip}{6pt plus 2pt minus 1pt}}
}{% if KOMA class
  \KOMAoptions{parskip=half}}
\makeatother
\setlength{\emergencystretch}{3em} % prevent overfull lines
\providecommand{\tightlist}{%
  \setlength{\itemsep}{0pt}\setlength{\parskip}{0pt}}
% v3 header: single column, tight reference list, clickable links.
% Loaded via pandoc -H. Fixes for the six issues Havid found in the v2 PDF are
% mostly in the markdown generator; this file only carries LaTeX-side settings.
\usepackage{graphicx}
\usepackage{float}
\usepackage{caption}

% Notation typesetting. Havid asked for consistency of format (2026-09-08),
% so units and species are expressed with the packages the field uses rather
% than a hand-rolled $^2$ regex table.
%
% VERIFIED before adoption: both compile under tectonic, which auto-downloads
% mhchem.sty and chemgreek.sty from its bundle -- no manual TeX Live install.
% Verified at PDF span level too: \qty{1.3e-3}{\centi\meter\squared\per\volt
% \per\second} raises every exponent AND every minus, and \ce{[PbI4]^2-} puts
% the subscript 4 BELOW the baseline with the charge ABOVE it in one token.
% A regex table cannot script one token in two directions.
%
% mhchem version=4 pinned: version 3 changed the charge syntax, so an
% unpinned load would silently alter every ion in the paper.
\usepackage{siunitx}
\usepackage[version=4]{mhchem}

% Journal-style units: upright glyphs, thin space between value and unit,
% and a real minus in exponents. detect-weight keeps units upright inside a
% bold heading instead of inheriting the surrounding bold.
\sisetup{
  detect-all,
  detect-weight        = true,
  range-units          = single,
  % reciprocal, NOT symbol. With per-mode=symbol the electron mobility
  % rendered as "cm2/(V s)", which is typographically correct but is not the
  % convention this literature uses; every perovskite paper writes
  % "cm2 V-1 s-1" with reciprocal powers. Verified at span level: reciprocal
  % raises both the 2 and each minus sign.
  per-mode             = reciprocal,
  exponent-product     = \times,
  output-decimal-marker = {.},
  group-digits         = none
}

% Figure captions: LaTeX owns the "Figure N." label. The v2 PDF rendered
% "Figure 4: Figure 4." because the caption text ALSO began with "Figure 4.",
% written into the markdown by the generator. v3 passes a bare caption and lets
% LaTeX number it, so the label appears exactly once.
\captionsetup[figure]{labelfont=bf, labelsep=period, font=small,
                      justification=justified, singlelinecheck=false,
                      skip=6pt}

% Tighten the numbered reference list produced by pandoc's default
% CSLReferences environment.
\usepackage{etoolbox}
\AtBeginEnvironment{CSLReferences}{%
  \small
  \setlength{\parskip}{0pt}%
  \setlength{\itemsep}{0pt}%
  \setlength{\topsep}{2pt}%
}

% No section numbering from LaTeX: the generator already writes "1. Introduction",
% "2. The Certified Frontier ...", so LaTeX numbering would double it.
\setcounter{secnumdepth}{0}
\usepackage{bookmark}
\IfFileExists{xurl.sty}{\usepackage{xurl}}{} % add URL line breaks if available
\urlstyle{same}
% fallback for those not using the hyperref driver hyperxmp:
\makeatletter
\@ifundefined{xmpquote}{\newcommand{\xmpquote}[1]{#1}}{}
\makeatother
\hypersetup{
  colorlinks=true,
  linkcolor={[HTML]{1A4E8A}},
  filecolor={Maroon},
  citecolor={Blue},
  urlcolor={[HTML]{1A4E8A}},
  pdfcreator={LaTeX via pandoc}}

\author{}
\date{}

\begin{document}

\setstretch{1.05}
\begin{center}
{\LARGE\bfseries Perovskite Photovoltaics in July 2026: Trap density and ion migration, self-assembled monolayer contacts\par}
\vspace{1.1em}
{\large Havid Aqoma\textsuperscript{1,2,3,4,*}\par}
\vspace{0.7em}
\begin{minipage}{\textwidth}\centering\small
\textsuperscript{1}School of Energy and Chemical Engineering, Xiamen University Malaysia, Selangor Darul Ehsan 43900, Malaysia\\[0.25em]
\textsuperscript{2}Kelip-kelip! Center of Excellence for Light Enabling Technologies, School of Energy and Chemical Engineering, Xiamen University Malaysia, Selangor Darul Ehsan 43900, Malaysia\\[0.25em]
\textsuperscript{3}Centre of Excellence for Green and Advanced Technologies in Efficient Separations (GATES), School of Energy and Chemical Engineering, Xiamen University Malaysia, Selangor Darul Ehsan 43900, Malaysia\\[0.25em]
\textsuperscript{4}Gulei Innovation Institute, Xiamen University, Zhangzhou 363200, China\\[0.55em]
*Corresponding author. Email: havidaqoma.khoiruddin@xmu.edu.my\\[2.0em]
September 9, 2026
\end{minipage}
\end{center}
\vspace{0.8em}

\subsection{Abstract}\label{abstract}

Across 648 indexed publications evaluated this month, with 181 examined
in depth, the verified performance ceiling reaches a power conversion
efficiency (PCE) of 33.1\% in monolithic tandems and 27.31\% in
single-junction inverted devices. Tandem benchmarks rely on current
matching across complementary bandgaps, whereas single-junction records
depend on multidentate chelating agents and self-assembled monolayer
(SAM) contacts that passivate undercoordinated lead cations and align
interfacial dipoles. Yet only 22 closely read studies provide
independently certified efficiency values, as uncertified reports
frequently mask mobile ionic defects and non-radiative recombination
under transient voltage sweeps. Scaling introduces severe resistive and
geometric transport losses, establishing an area penalty that restricts
certified efficiency to 27.31\% at the largest reported aperture area of
\qty{766}{\centi\meter}\(^2\). This scaling deficit arises because
molecular passivators that eliminate trap states on microscale cells
aggregate during meniscus drying, intensifying non-radiative
recombination across laser-scribed module interconnects. Operational
durability evidence remains narrow across the literature: only 4
investigations quantify lifetime metrics, with the longest duration
reaching \qty{9800}{\hour}, while only 6 studies adhere to standardized
International Summit on Organic Photovoltaic Stability (ISOS) protocols.
A pronounced reporting disparity persists between claimed and verified
performance, with 50.4\% of abstracts reporting efficiency but only
8.3\% documenting independent certification. Mechanistic evaluation
demonstrates that contact stabilization through self-assembled
monolayers (SAMs) and multidentate coordination fails to establish
resilience against thermo-mechanical fatigue, as thermal expansion
mismatches induce cyclic shear stresses that trigger interfacial
delamination. Depth-resolved degradation profiles confirm that breakdown
initiates at buried selective contacts before penetrating bulk
absorbers. Chemical treatments and interlayers also fail to achieve
permanent suppression of halide phase segregation in wide-bandgap
absorbers, where illumination and thermal cycling drive phase demixing
without resolving the boundary between thermodynamic stability and
kinetic demixing arrest. Finally, unstandardized aging without
continuous tracking leaves the division between reversible ionic
polarization and irreversible lattice degradation unquantified.

\textbf{Keywords:} perovskite photovoltaics; buried interface; defect
passivation; wide-bandgap perovskite; tandem solar cells; operational
stability

\subsection{1. Introduction}\label{introduction}

Perovskite photovoltaics has reached a juncture where incremental gains
in laboratory power conversion efficiency (PCE) no longer mark material
progress. The defining frontier has shifted to independent verification,
scalability across module-relevant area, and operating lifetime under
operational conditions. Peak metrics obtained from uncertified,
sub-millimetre devices under transient voltage sweeps frequently obscure
mobile ionic defects and non-radiative interfacial recombination that
manifest only over extended timescales.

Within this landscape, 648 indexed works met the scope gate in July 2026
and 181 were read closely.

The mechanistic threads followed in this issue encompass passivation
chemistry and ion migration under bias, self-assembled contacts and
interfacial dipoles, precursor chemistry and phase purity, stability
under load, heat and illumination, and wide-bandgap subcells and
recombination layers. Together, these processes govern whether energetic
band alignment and chemical defect suppression survive the coupled
electrochemical gradients of field operation.

To evaluate these developments, this review distinguishes what the
month's literature verified from what it asserted. Filtering out
artifacts of fast scan rates, uncalibrated spectral illumination, and
negligible active areas isolates the specific chemical and structural
interventions that provide authentic stability under working conditions.

\subsection{2. The Certified Frontier in July
2026}\label{the-certified-frontier-in-july-2026}

The verified performance frontier reaches a power conversion efficiency
(PCE) of 33.1\% in two-terminal (2T) tandem architectures and 27.31\% in
single-junction inverted (p-i-n) devices. Yet only 22 of 169 closely
read papers reported an independently certified efficiency. In 2T
tandems, certified values of 33.1\% and 32.13\%, along with intermediate
benchmarks at 29.57\%, 28.84\%, and 27.49\%, confirm that current
matching across paired subcells dictates the verifiable ceiling
{[}\href{https://doi.org/10.1126/sciadv.aei2945}{1},\href{https://doi.org/10.1002/anie.1676645}{2},\href{https://doi.org/10.1002/adma.74078}{3},\href{https://doi.org/10.1126/sciadv.aeb8790}{4},\href{https://doi.org/10.1126/sciadv.aef6600}{5}{]}.
Across p-i-n single junctions, certified PCE concentrates tightly at
27.31\% and 27.3\%, establishing an interfacial recombination barrier
where non-radiative losses at transport layers suppress open-circuit
voltage
{[}\href{https://doi.org/10.1002/adma.73807}{6},\href{https://doi.org/10.1126/sciadv.aeg1456}{7},\href{https://doi.org/10.1126/science.aed8175}{8}{]}.

Area expansion reveals how uncertified laboratory reports outpace
verified metrics. For a \qty{625}{\centi\meter\squared} aperture,
self-reported efficiency reached 27.8\% despite an independently
certified baseline of 27.3\%
{[}\href{https://doi.org/10.1126/science.aed8175}{8}{]}. Certified
large-area devices quantify the resistive and geometric penalties
imposed by macroscopic transport. An \qty{818}{\centi\meter\squared}
device achieved 26.53\%
{[}\href{https://doi.org/10.1021/jacs.6c06932}{9}{]}. A
\qty{766}{\centi\meter\squared} device sustained 27.31\%
{[}\href{https://doi.org/10.1126/sciadv.aeg1456}{7}{]}. A
\qty{684.75}{\centi\meter\squared} device reached 27.31\%
{[}\href{https://doi.org/10.1002/adma.73807}{6}{]}. A
\qty{655.2}{\centi\meter\squared} device delivered 27.12\%
{[}\href{https://doi.org/10.1038/s41467-026-74018-8}{10}{]}. These
metrics verify that scalable interconnects and spatial uniformity can
suppress series resistance across large modules.

A champion efficiency obtained from an unshadowed,
sub-\qty{0.1}{\centi\meter\squared} laboratory cell and an independently
certified measurement across a module represent irreconcilable
categories of physical evidence. Microscale test structures probe
localized carrier generation under negligible series resistance,
systematically masking pinhole shunting, film non-uniformity, and
perimeter recombination. Conversely, certified module-scale metrics
integrate transparent electrode sheet resistance, scribe dead zones, and
macroscopic phase segregation across active layers. Equating isolated
laboratory spikes with certified module performance mistakes local
material potential for verified device technology.

\begin{figure}[htbp]
\centering
\includegraphics[width=0.95\textwidth]{fig/F1_certified_frontier.pdf}
\caption{The certified frontier in July 2026. Filled markers are independently certified efficiencies, open markers are self-reported champion values, grouped by device family. Horizontal bars mark the highest certified value in each family.}
\end{figure}

\subsection{3. Passivation Chemistry and Ion Migration Under
Bias}\label{passivation-chemistry-and-ion-migration-under-bias}

The quantitative frontier for perovskite solar cells now reaches
certified power conversion efficiency (PCE) values of 26.84\% (champion
27.25\%) in inverted devices and 23.5\% on
\qty{65.05}{\centi\meter\squared} modules (champion 27.2\%). Uncertified
small-area claims reach 26.85\%
{[}\href{https://doi.org/10.1002/aenm.71305}{11},\href{https://doi.org/10.1126/science.aeg1730}{12},\href{https://doi.org/10.1002/anie.1864800}{13}{]}.
These benchmarks force a critical evaluation of what the reported
physical evidence actually licenses. One possibility is that interfacial
chemical treatments permanently suppress deep trap states and mobile ion
channels under operating bias. The alternative is that they merely
provide transient electrostatic screening that degrades during
continuous operation.

At perovskite interfaces, undercoordinated \ce{Pb^2+} cations and halide
vacancies constitute primary non-radiative recombination centers.
Monodentate Lewis bases donate electron pairs from sulfur, nitrogen, or
oxygen to passivate surface \ce{Pb^2+} defects. Carbazole-based
self-assembled monolayers (SAMs) such as BTP at buried interfaces,
thiophene sulfur coordination, and carboxyl/sulfur bifunctional TTA
demonstrate this route. These three report uncertified champion PCE
values of 26.85\%, 21.15\%, and 26.7\%
{[}\href{https://doi.org/10.1002/anie.1864800}{13},\href{https://doi.org/10.1021/acsami.6c05331}{14},\href{https://doi.org/10.1002/adma.73950}{15}{]}.
However, monodentate coordination offers limited thermal resistance
against cooperative defect formation. To address this, amphoteric ionic
liquids (GUHS and GUDHP) and FTAC simultaneously passivate Lewis acidic
and basic sites while compensating iodine vacancies. GUDHP achieves
20.06\% PCE with 85.74\% fill factor (FF) at \qty{1.84}{\electronvolt}
bandgap (champion 26.13\%). FTAC retains 88.67\% of initial efficiency
after 1,\qty{000}{\hour} maximum-power-point tracking (maximum power
point tracking (MPPT)) at 65 °C despite an uncertified 25.06\% champion
PCE
{[}\href{https://doi.org/10.1002/adfm.77378}{16},\href{https://doi.org/10.20517/energymater.2026.102}{17}{]}.
Multivalent spatial anchoring further impedes defect formation.
Three-dimensional distribution of amino groups in
tetrakis(4-aminophenyl)methane (4APM), nonlayered intermediate-site
anchoring ligands, and m-bis(trifluoromethyl)aniline (mTFMA) bind
\ce{PbI2} or formamidinium iodide (FAI) precursors. That binding
lengthens carrier lifetimes from 529.7 to \qty{2479.4}{\nano\second}.
These films retain over 95\% efficiency after \qty{1000}{\hour} MPPT at
65 °C across \qty{14}{\centi\meter\squared} (certified 26.26\%). They
also retain over 90\% PCE after \qty{2000}{\hour} at 65 °C and 90\%
relative humidity (RH) (uncertified champion 26.54\%, projected T80 of
\qty{9800}{\hour})
{[}\href{https://doi.org/10.1038/s41467-026-75990-x}{18},\href{https://doi.org/10.1002/adma.74243}{19},\href{https://doi.org/10.1021/acsaem.6c01560}{20}{]}.

Under continuous bias, applied electric fields drive mobile iodide
migration toward extraction layers, causing interfacial degradation and
phase segregation. Chemical immobilization directly arrests this
migration. Dual iodide-trapping agents (TFIBA and FPIFA),
two-dimensional/three-dimensional (2D/3D) halide imidazolium
heterojunctions, and 1-(4'-sulfophenyl)-3-carboxyl-5-pyrazolone sodium
salt (SCP) layers suppress spiro-OMeTAD and electron-contact
degradation. The first route achieves an uncertified 26.81\% champion
PCE with T93 lifetime over \qty{1140}{\hour} at 85 °C. The second holds
92\% retention after \qty{1200}{\hour} at 85 °C (champion 26.33\%). The
third holds 90\% retention after \qty{1485}{\hour} MPPT (certified
26.84\%, champion 27.25\%)
{[}\href{https://doi.org/10.1002/aenm.71305}{11},\href{https://doi.org/10.1002/anie.2924886}{21},\href{https://doi.org/10.1002/adfm.77015}{22}{]}.
In wide-bandgap absorbers, field-driven halide separation is suppressed
by guanidine phosphate (champion 21.54\% at \qty{1.78}{\electronvolt},
retaining 90.64\% after \qty{1400}{\hour}). It is also suppressed by
acetylthiocholine chloride (ATC) doping, which fills vacancies with
quaternary ammonium and chloride ions to restrain demixing across
\qty{63.7}{\centi\meter\squared} modules (champion 22.68\%)
{[}\href{https://doi.org/10.1002/cnl2.70189}{23},\href{https://doi.org/10.1002/ece2.70102}{24}{]}.
Conformal plasma surface engineering similarly mitigates defect
accumulation at laser-scribed interconnects to secure a certified 23.5\%
module PCE on \qty{65.05}{\centi\meter\squared} (champion 27.2\%).
Small-area devices retained 98.1\% of initial PCE after
\qty{2000}{\hour} MPPT at 85 °C
{[}\href{https://doi.org/10.1126/science.aeg1730}{12}{]}.

Lattice strain fundamentally governs ion migration kinetics by lowering
vacancy hopping barriers under tensile expansion. Relieving this strain
through dual-site coordination with dithiopyr bridges adjacent
\ce{Pb^2+} sites, promoting preferential (100) crystal orientation and
yielding a certified 26.14\% PCE (champion 26.9\%, and 21.28\% on 30 ×
\qty{30}{\centi\meter\squared} modules)
{[}\href{https://doi.org/10.1002/adma.74033}{25}{]}. Compressive strain
engineering counteracts tensile distortion and impedes defect formation
through three routes. The polymer MPUA reaches a champion 24.91\%.
Cross-linked 2-amino-3-hydroxybutanoic acid (AHBA) interfaces accelerate
electron extraction to \qty{3.5}{\pico\second} (champion 26.32\%).
Photoresponsive \ce{C60}-azo trans-to-cis isomerization under
ultraviolet (UV) light gives 92.7\% retention after \qty{488}{\hour}
(champion 26.6\%)
{[}\href{https://doi.org/10.1002/solr.70435}{26},\href{https://doi.org/10.1002/adfm.77373}{27},\href{https://doi.org/10.1002/adma.74220}{28}{]}.
Nevertheless, most strain and passivation studies report uncertified,
small-area metrics without operando tracking of mobile ion accumulation
under concurrent bias, heat, and illumination, leaving long-term field
immunity unsubstantiated.

\begin{figure}[htbp]
\centering
\includegraphics[width=0.95\textwidth]{fig/F3_effort_map.pdf}
\caption{Where the month's mechanistic effort was directed. Depth-tier papers by mechanism axis, stacked by device family.}
\end{figure}

\subsection{4. Self-Assembled Contacts and Interfacial
Dipoles}\label{self-assembled-contacts-and-interfacial-dipoles}

The quantitative frontier for contact-engineered perovskite
photovoltaics is established by monolithic tandem configurations
achieving certified power conversion efficiency (PCE) values of 33.1\%
(champion 33.3\%) {[}\href{https://doi.org/10.1126/sciadv.aei2945}{1}{]}
and single-junction inverted architectures reaching certified 27.31\%
{[}\href{https://doi.org/10.1002/adma.73807}{6}{]}. Translating these
efficiencies into durable devices hinges on resolving three core
physical problems: weak oxide anchoring, self-assembled monolayer (SAM)
aggregation, and interfacial dipole misalignment.

Substrate delamination and thermal desorption originate from fragile
monodentate coordination at transparent conducting oxides. Mechanical
toughness is markedly substrate-dependent: depositing self-assembled
monolayers (SAMs) on a nickel oxide (\ce{NiO_x}) underlayer increases
interfacial fracture energy tenfold relative to indium tin oxide (ITO),
while removing loosely bound molecules elevates fracture energy from
0.94 to \qty{4.82}{\joule\per\meter\squared}
{[}\href{https://doi.org/10.1021/acsmaterialslett.6c00479}{29}{]}.
Multidentate chelating motifs directly counter this vulnerability by
coordinating undercoordinated lead (\ce{Pb^2+}) and reinforcing
substrate anchoring. Tetradentate iron(III) porphyrins yield certified
26.82\% PCE (champion 27.06\%)
{[}\href{https://doi.org/10.1002/anie.8080769}{30}{]}. Bisphosphonates
reconstruct the \ce{NiO_x}/SAM buried interface to reach certified
27.31\% PCE on small-area rigid cells
{[}\href{https://doi.org/10.1002/adma.73807}{6}{]}. Zoledronic acid
(ZDA) head groups yield 30.62\% tandem PCE, and the wide-bandgap
sub-cells (20.55\% PCE) retain over 90\% performance after
\qty{200}{\hour} of maximum power point (MPP) tracking
{[}\href{https://doi.org/10.1002/adma.73807}{6},\href{https://doi.org/10.1002/anie.8080769}{30},\href{https://doi.org/10.1002/smll.74999}{31}{]}.
Lindqvist-type polyoxometalates similarly raise the anchoring ratio on
ITO from 29.93\% to 53.98\% and cut the desorption rate from 30.7\% to
5.6\% during continuous MPP tracking (champion 25.15\% PCE)
{[}\href{https://doi.org/10.26599/pom.2026.9140147}{32}{]}. Large-area
module demonstrations remain scarce, though tetraphosphorylated
phthalocyanines demonstrate a certified 23.1\% PCE (stabilised 22.9\%)
on a \qty{20.3}{\centi\meter\squared} module without degradation after
61 days of International Summit on Organic Photovoltaic Stability
(ISOS)-O-1 outdoor aging
{[}\href{https://doi.org/10.1126/sciadv.aee4175}{33}{]}.

Intermolecular pi-pi stacking among aromatic cores drives nanoscale SAM
agglomeration, generating pinholes and phase inhomogeneities. Disrupting
these interactions via steric hindrance or electrostatic confinement
suppresses clustering. Noncoplanar multidentate contacts (TTA) impede
pi-pi stacking to stabilize wide-bandgap single-junction absorbers
(approximately \qty{1.67}{\electronvolt}, champion PCE of 23.7\%) and
tandems {[}\href{https://doi.org/10.1126/sciadv.aei2945}{1}{]}. Ionic
liquid lubricants (BMIMBF4) suppress aggregation via dual
steric-electrostatic modulation. Those devices retained 92.3\% of
initial efficiency after \qty{800}{\hour} at 85 °C under air mass 1.5
global (AM 1.5G) illumination (uncertified champion 26.82\% PCE)
{[}\href{https://doi.org/10.1002/anie.7090517}{34}{]}. Concurrently,
proton transfer complex-confined carbazole SAMs maintain thermal
stability at 75 °C for \qty{1368}{\hour} (champion 26.23\% PCE)
{[}\href{https://doi.org/10.1021/acsenergylett.6c01489}{35}{]}, and in
situ thermally activated polymerization networks seal interfacial voids
to achieve 26.87\% champion PCE
{[}\href{https://doi.org/10.1002/anie.7090517}{34},\href{https://doi.org/10.1021/acsenergylett.6c01489}{35},\href{https://doi.org/10.1002/anie.1018855}{36}{]}.

Unfavorable vacuum-level shifts and unpassivated defects dictate carrier
extraction deficits at the absorber boundary. Halogenation and
co-assembly modulate work functions. Substituting methoxy groups with
fluorine in F-2PACz strengthens the interfacial dipole and passivates
halide vacancies. That contact retains 25.06\% PCE on
\qty{1}{\centi\meter\squared} devices with 3.11\% relative efficiency
loss (champion 25.75\%)
{[}\href{https://doi.org/10.1002/adfm.76878}{37}{]}. Co-assembling crown
ethers at electron transport layer interfaces instead increases surface
coverage from 0.57 to 0.79 and strengthens interfacial dipole moments
from 2 to 7 Debye (champion 26.33\% PCE)
{[}\href{https://doi.org/10.1002/adma.73886}{38}{]}. Polycyclic
perylenocarbazole cores (6PCzBr) tailor electrostatic potentials for
26.20\% PCE {[}\href{https://doi.org/10.1002/anie.4031162}{39}{]}.
Phosphine additives (TFMP) coordinate \ce{Pb^2+} and stabilize
formamidinium (FA+) in Me-4PACz and 4PABCz ternary contacts (champion
26.78\% PCE) {[}\href{https://doi.org/10.1002/adma.74416}{40}{]}.
Ammonium bromide interlayers (PFN-Br) tune band alignment for 23.15\%
PCE {[}\href{https://doi.org/10.1002/smll.74806}{41}{]}. However,
because most dipolar gains are reported on uncertified
sub-\qty{0.1}{\centi\meter\squared} cells without standardized bias
protocols, whether these electrostatic barriers resist operational ion
drift remains unverified.

\subsection{5. Precursor Chemistry and Phase
Purity}\label{precursor-chemistry-and-phase-purity}

The quantitative frontier for perovskite photovoltaics is defined by
single-junction power conversion efficiency (PCE) reaching certified
values of 26.52\% (26.88\% champion)
{[}\href{https://doi.org/10.1126/sciadv.aed6327}{42}{]} and uncertified
reports exceeding 26\%
{[}\href{https://doi.org/10.1002/advs.76508}{43},\href{https://doi.org/10.26434/chemrxiv.15006632/v1}{44}{]}.
Whether these metrics reflect phase-pure lattices or metastable
polymorphic mixtures remains governed by precursor kinetics, phase
segregation, and spatial scalability. Record efficiencies frequently
rely on sub-\qty{0.1}{\centi\meter\squared} apertures, such as 26.32\%
on \qty{0.072}{\centi\meter\squared}
{[}\href{https://doi.org/10.1002/advs.76508}{43}{]}. Scaling to a
\qty{900}{\centi\meter\squared} module lowers efficiency to 20.25\%
(17.75\% certified). That module still retained over 91\% initial
efficiency after 1000 hours at 85°C and 85\% relative humidity (RH)
{[}\href{https://doi.org/10.1126/sciadv.aed6327}{42}{]}.

Controlling precursor coordination prevents solvent degradation and
defective phase precipitation
{[}\href{https://doi.org/10.1021/jacs.6c05502}{45},\href{https://doi.org/10.1021/acsaem.6c01764}{46},\href{https://doi.org/10.1002/adfm.76731}{47}{]}.
In formamidinium (FA) lead triiodide (\ce{FAPbI3}) growth, parasitic
esterification between gamma-butyrolactone (GBL) and 2-methoxyethanol
(2ME) destabilizes heated solutions; maintaining lead iodide (\ce{PbI2})
deficiency delays redissolution to suppress delta-phase impurities,
yielding 22.99\% PCE
{[}\href{https://doi.org/10.1021/jacs.6c05502}{45}{]}. Additives also
eliminate antisolvent quenching: sulfolane forms an intermediate phase
yielding 22.2\% PCE in methylammonium (MA) lead triiodide (\ce{MAPbI3})
{[}\href{https://doi.org/10.1021/acsaem.6c01764}{46}{]}, while 3-furoic
acid (3-FA) coordinates tin iodide (\ce{SnI2}) and hydrogen-bonds with
formamidinium iodide (FAI) to decouple nucleation from crystal growth,
reaching 16.16\% PCE
{[}\href{https://doi.org/10.1002/adfm.76731}{47}{]}. Stoichiometric
imbalance leaves defective binary halides: PbX2-terminated surfaces
display higher defect densities and lower stability than AX terminations
{[}\href{https://doi.org/10.1021/acsnano.6c08922}{48}{]}. Passivating
residual \ce{PbI2} with triphenyl phosphoramide (TPPA) yields 21.4\% PCE
and retains 90\% initial efficiency after 750 hours of maximum power
point tracking (MPPT) versus 270 hours for controls
{[}\href{https://doi.org/10.1021/acsaem.6c01370}{49}{]}. Similarly,
molten 4-(4,6-dimethoxy-1,3,5-triazin-2-yl)-4-methylmorpholinium
chloride (DMTMM) dissolves two-dimensional (2D) phases to eliminate
interfacial \ce{PbI2}, raising PCE from 23.43\% to 25.73\%
{[}\href{https://doi.org/10.1002/adfm.76974}{50}{]}, while optimizing
\ce{PbI2} evaporation distance achieves 14.57\% PCE over
\qty{64}{\centi\meter\squared}
{[}\href{https://doi.org/10.1002/solr.70430}{51}{]}.

Beyond nucleation, compositional heterogeneity drives phase demixing
{[}\href{https://doi.org/10.1002/adma.74382}{52},\href{https://doi.org/10.1016/j.solmat.2026.114597}{53}{]}.
Concurrent four-dimensional scanning transmission electron microscopy
(4D-STEM) and energy-dispersive X-ray spectroscopy (EDX) reveals that
photoinactive hexagonal polytypes nucleate from nanoscale cesium (Cs)
deficiencies {[}\href{https://doi.org/10.1002/adma.74382}{52}{]}. In
mixed-halide absorbers, strontium doping prevents I2Br- intermediate
formation, suppressing segregation to reach 20.18\% PCE at
\qty{1.77}{\electronvolt} and 19.02\% at \qty{1.84}{\electronvolt}
{[}\href{https://doi.org/10.1021/acsenergylett.6c01113}{54}{]}. Rubidium
A-site incorporation induces lattice distortion to suppress halide
demixing, providing an open-circuit voltage (Voc) of \qty{1.46}{\volt}
in wide-bandgap sub-cells and \qty{3.31}{\volt} in triple-junction
tandems {[}\href{https://doi.org/10.1002/adfm.76715}{55}{]}. Conversely,
trivalent bismuth substitution in cesium lead halides (CsPbX3) triggers
phase segregation due to lower thermodynamic stability
{[}\href{https://doi.org/10.1016/j.solmat.2026.114597}{53}{]}.
Dielectric-response screening of 5043 compositions confirms that phase
stability requires narrow windows of 2--8\% Cs, 78--89\% FA, 3--18\% MA,
and 93--99\% iodide (I)
{[}\href{https://doi.org/10.26434/chemrxiv.15005541/v1}{56}{]}.

Dimensional heterostructuring stabilizes interfaces, but phase purity
remains constrained by transport penalties and area scaling
{[}\href{https://doi.org/10.1002/advs.76508}{43},\href{https://doi.org/10.26434/chemrxiv.15006632/v1}{44},\href{https://doi.org/10.1002/anie.7827890}{57}{]}.
In situ hydrogen-bonded hydroxyphenethylammonium (HPA) spacer networks
retain 95.5\% of initial efficiency after 3000 hours at 30°C in \ce{N2},
92.5\% in 45\% ± 5\% RH air, and 90\% after 800 hours at 85°C in
\ce{N2}, but on \qty{0.072}{\centi\meter\squared} cells at 26.32\% PCE
{[}\href{https://doi.org/10.1002/advs.76508}{43}{]}. Colloidal
suspension epitaxy yields deterministic 26.0\% multilayers without
disturbing underlying layers
{[}\href{https://doi.org/10.26434/chemrxiv.15006632/v1}{44}{]}. In
flexible tin cells, interfacial hydrogen bonding yields 15.56\% PCE
(14.67\% certified) on \qty{0.049}{\centi\meter\squared}, but drops to
11.06\% at \qty{1}{\centi\meter\squared} despite achieving 22.21\% under
1000 lux indoor illumination
{[}\href{https://doi.org/10.1002/anie.7827890}{57}{]}. While precursor
coordination and dimensional stabilization preserve microscopic phase
purity, sub-\qty{0.1}{\centi\meter\squared} testbeds obscure severe
interfacial penalties encountered during physical scaling.

\subsection{6. Stability Under Load, Heat and
Illumination}\label{stability-under-load-heat-and-illumination}

The quantitative frontier for operational durability is defined by time
to 80\% initial performance (T80) values of \qty{9800}{\hour}
{[}\href{https://doi.org/10.1002/adma.74243}{19}{]} and
\qty{6785}{\hour} {[}\href{https://doi.org/10.1002/aenm.71317}{58}{]},
falling steeply to \qty{720}{\hour}
{[}\href{https://doi.org/10.1021/acsaem.6c01637}{59}{]} and
\qty{200}{\hour} {[}\href{https://doi.org/10.1002/smll.74531}{60}{]}.
However, this frontier rests on narrow evidence. Only 4 of 169 closely
read studies report a T80 metric. Only 6 identify standardized
International Summit on Organic Photovoltaic Stability (ISOS) protocols,
such as ISOS-O-1
{[}\href{https://doi.org/10.1126/sciadv.aee4175}{33}{]}, ISOS-L-1
{[}\href{https://doi.org/10.1126/sciadv.aeg1456}{7},\href{https://doi.org/10.1021/acsnano.6c04084}{61}{]},
ISOS-L-3 {[}\href{https://doi.org/10.1021/acsenergylett.6c01517}{62}{]},
ISOS-L-1l {[}\href{https://doi.org/10.1021/jacs.6c06932}{9}{]}, or
continuous maximum power point tracking (MPPT)
{[}\href{https://doi.org/10.26599/pom.2026.9140147}{32}{]}. Most
stability claims instead rely on unstandardized, short-duration testing.

Depth-resolved profiling reveals that photo-induced degradation
originates at both top and buried interfaces before penetrating the bulk
{[}\href{https://doi.org/10.1002/solr.70421}{63}{]}. To arrest
interfacial chemical breakdown, molecularly tailored self-assembled
monolayers (SAMs) stabilize buried contacts. These contacts yield a
certified power conversion efficiency (PCE) of 27.31\% in single cells
and 21.54\% on a \qty{766}{\centi\meter\squared} module under ISOS-L-1
{[}\href{https://doi.org/10.1126/sciadv.aeg1456}{7}{]}. Asymmetric SAMs
achieve 20.59\% on \qty{655.2}{\centi\meter\squared} modules without
certified durability
{[}\href{https://doi.org/10.1002/anie.1678852}{64}{]}. At the upper
contact, multidentate fullerenes passivate defects and immobilize
surface cations under ISOS-L-3, retaining 98.7\% of initial PCE after
\qty{2100}{\hour} of continuous illumination and 95.8\% after
\qty{1800}{\hour} at 85 °C, though on uncertified devices
{[}\href{https://doi.org/10.1021/acsenergylett.6c01517}{62}{]}.
Interfacial defect passivation via aromatic carbonyls similarly claims
stability but remains confined to an uncertified,
sub-\qty{0.1}{\centi\meter\squared} cell of
\qty{0.09}{\centi\meter\squared}
{[}\href{https://doi.org/10.1002/anie.5216345}{65}{]}. In-situ
two-dimensional perovskite reconstruction further blocks interfacial
degradation
{[}\href{https://doi.org/10.1002/adma.74247}{66},\href{https://doi.org/10.1002/solr.70447}{67}{]},
while volatile dopants eliminate mobile lithium-ion residues in
hole-transporting layers
{[}\href{https://doi.org/10.1002/smll.74289}{68}{]}.

Under continuous bias and illumination, trapped-charge accumulation
accelerates halide phase segregation, degrading wide-bandgap absorbers
under standard MPPT, whereas maximum stability point tracking (MSPT)
retains approximately 90\% of initial power after \qty{700}{\hour}
{[}\href{https://doi.org/10.1002/aenm.71358}{69}{]}. Ion-migration
barriers from diammonium additives suppress segregation to achieve a
certified 26.18\% PCE on a \qty{1.17}{\centi\meter\squared}
triple-junction aperture
{[}\href{https://doi.org/10.1038/s41467-026-75305-0}{70}{]}, aligning
with rubidium doping to homogenize vertical cation distributions under
ISOS-L-1 {[}\href{https://doi.org/10.1021/acsnano.6c04084}{61}{]}. These
degradation pathways are intensified by dynamic thermo-mechanical stress
fluctuations
{[}\href{https://doi.org/10.1021/acsenergylett.6c01837}{71}{]}, where
inelastic tensile stress pre-relaxation reduces moisture ingress by up
to an order of magnitude
{[}\href{https://doi.org/10.26434/chemrxiv-2025-128tn/v2}{72}{]}.
Without widespread adoption of standardized ISOS protocols, reported
operational lifetimes remain poorly substantiated.

\begin{figure}[htbp]
\centering
\includegraphics[width=0.95\textwidth]{fig/F4_stability_evidence.pdf}
\caption{The month's operational-stability evidence in full. Left, every reported time to 80\% of initial performance. Right, the stability protocol labels used.}
\end{figure}

\subsection{7. Wide-Bandgap Subcells and Recombination
Layers}\label{wide-bandgap-subcells-and-recombination-layers}

Monolithic perovskite-silicon tandems have pushed power conversion
efficiency (PCE) to 33.21\% (certified 32.13\%) and 32.5\% through
interface and crystallization engineering
{[}\href{https://doi.org/10.1002/anie.1676645}{2},\href{https://doi.org/10.1007/s40820-026-02291-9}{73}{]},
but reported performance exceeds what operating stability currently
licenses. Two physical impediments dominate: interfacial recombination
at the hole-selective contact and halide phase segregation in
wide-bandgap (WBG) absorbers.

At the buried contact, self-assembled monolayers (SAMs) suffer from
solvent-induced degradation and inhomogeneous nucleation. Inserting a
3-TMA molecular interlayer between the self-assembled monolayer (SAM)
and the absorber regulates crystallization via hydrogen bonding. That
interlayer reaches 23.17\% PCE in \qty{1.67}{\electronvolt} and 19.61\%
in \qty{1.84}{\electronvolt} single-junction devices
{[}\href{https://doi.org/10.1002/anie.1676645}{2}{]}. Combining graphene
oxide with SAMs forms a contact that retains 90\% of initial efficiency
after \qty{770}{\hour} at 25 °C in a 27.3\% tandem
{[}\href{https://doi.org/10.1016/j.joule.2026.102575}{74}{]}.
Coordinating formamidinium cations with gamma-valerolactone (GVL)
cosolvents similarly retards nucleation to enforce oriented
micrometer-scale grain growth across the buried interface
{[}\href{https://doi.org/10.1007/s40820-026-02291-9}{73},\href{https://doi.org/10.1002/advs.76463}{75}{]}.

Within the bulk, mixed-halide absorbers undergo light- and heat-induced
phase segregation. Under in-operando thermal cycling from -100 to +100
°C, WBG subcells exhibit severe phase segregation above 50 °C with only
partial recovery during dark cooling
{[}\href{https://doi.org/10.1021/acsenergylett.6c00992}{76}{]}. Halogen
homogenization via TPB-Br surface passivation
{[}\href{https://doi.org/10.1002/anie.5530609}{77}{]} and conductive
polymer grain-boundary bridges
{[}\href{https://doi.org/10.1002/adma.74183}{78}{]} suppress halide
demixing. These routes yield 23.63\% (\qty{1.68}{\electronvolt} subcell)
and 20.85\% single junctions, alongside uncertified tandems of 32.03\%
on \qty{1.05}{\centi\meter\squared} and 25.92\%
{[}\href{https://doi.org/10.1002/anie.5530609}{77},\href{https://doi.org/10.1002/adma.74183}{78}{]}.
Strengthening the metal-halide octahedral framework via \ce{Sr^2+}
alloying weakens electron-phonon coupling, achieving 31.66\% PCE with an
operational lifetime to 80\% of initial efficiency (T80) of
\qty{6785}{\hour} on fully-textured silicon
{[}\href{https://doi.org/10.1002/aenm.71317}{58}{]}.

Upscaling these stacks exposes severe area-dependent transport
penalties. Quantitative hyperspectral luminescence imaging on
\qty{64}{\centi\meter\squared} modules reveals that resistive transport
and interconnection losses equal intrinsic material losses
{[}\href{https://doi.org/10.1002/solr.70375}{79}{]}. Scalable printing
using low-toxicity GVL with dimethyl sulfoxide (DMSO) exceeds 17\% in
WBG cells and 28\% in tandems
{[}\href{https://doi.org/10.1002/advs.76463}{75}{]}, while ionic liquid
stabilization demonstrates certified 27.49\% tandems on
\qty{12.56}{\centi\meter\squared} with 23.58\% certified top-cell
efficiency {[}\href{https://doi.org/10.1126/sciadv.aef6600}{5}{]}.

\begin{figure}[htbp]
\centering
\includegraphics[width=0.95\textwidth]{fig/F2_area_penalty.pdf}
\caption{Efficiency against reported active or aperture area. The shaded band marks the sub-0.3 cm$^2$ laboratory-cell regime in which most champion values are measured.}
\end{figure}

\subsection{8. Research Gaps and
Outlook}\label{research-gaps-and-outlook}

The verified performance ceiling reaches a certified power conversion
efficiency (PCE) of 33.1\% in two-terminal (2T) tandems
{[}\href{https://doi.org/10.1126/sciadv.aei2945}{1}{]} and 27.31\% in
single junctions
{[}\href{https://doi.org/10.1002/adma.73807}{6},\href{https://doi.org/10.1126/sciadv.aeg1456}{7}{]}.
Yet the certified reporting share was 5.9\% in 2026-06 and 8.3\% in
2026-07, with only 22 of 169 closely read papers reporting any certified
value. These peak metrics do not license claims of operating durability,
as uncertified microscale sweeps mask interfacial degradation, phase
instability, and scaling penalties.

First, evidence for contact stabilization fails to establish whether
interfacial dipoles and multidentate coordination withstand
thermo-mechanical fatigue. While modifying transparent oxides with
self-assembled monolayer (SAM) contacts and multidentate chelators
elevates initial open-circuit voltage (Voc), thermal expansion
mismatches induce cyclic shear stresses that trigger delamination
{[}\href{https://doi.org/10.1126/sciadv.aei2945}{1},\href{https://doi.org/10.1002/adma.73807}{6},\href{https://doi.org/10.1021/acsenergylett.6c01489}{35}{]}.
Resolving whether coordination bonds survive operating fields requires
operando double-cantilever beam fracture testing paired with
bias-dependent Kelvin probe force microscopy during thermal cycling.

Second, reported chemical treatments do not demonstrate permanent
suppression of halide segregation in wide-bandgap (WBG) absorbers.
Although molecular interlayers and B-site alloying achieve high initial
efficiency in tandems
{[}\href{https://doi.org/10.1002/anie.1676645}{2},\href{https://doi.org/10.1002/aenm.71317}{58}{]},
in situ thermal cycling reveals phase segregation above 50 °C with only
partial recovery upon cooling
{[}\href{https://doi.org/10.1021/acsenergylett.6c00992}{76}{]}.
Differentiating true thermodynamic phase stability from kinetic demixing
arrest demands operando nanoscale photoluminescence spectroscopy
synchronized with diffraction profiling under continuous electrical bias
and elevated temperature.

Third, unstandardized testing leaves the boundary between reversible
ionic polarization and irreversible decomposition unquantified. Only 4
papers reported a time to 80\% of initial performance (T80) lifetime,
the longest \qty{9800}{\hour}
{[}\href{https://doi.org/10.1002/adma.74243}{19}{]}, and only 6 papers
named an International Summit on Organic Photovoltaic Stability (ISOS)
protocol
{[}\href{https://doi.org/10.1021/acsenergylett.6c01517}{62},\href{https://doi.org/10.1002/aenm.71358}{69}{]}.
Establishing ion accumulation rates under operating conditions requires
operando intensity-modulated photovoltage spectroscopy coupled with
secondary-ion mass spectrometry during maximum power point tracking
(MPPT).

Fourth, microscale defect passivation fails to establish scalability
across large apertures. Only 27 of 648 works addressed scale-up, and the
largest reported aperture reached \qty{818}{\centi\meter\squared}
{[}\href{https://doi.org/10.1021/jacs.6c06932}{9}{]}. Passivators that
eliminate traps on sub-\qty{0.1}{\centi\meter\squared} apertures
aggregate during meniscus drying, intensifying non-radiative
recombination across laser-scribed modules
{[}\href{https://doi.org/10.1126/sciadv.aed6327}{42},\href{https://doi.org/10.1002/solr.70375}{79}{]}.
Spatially resolved electroluminescence imaging calibrated against local
series-resistance maps during coating would isolate interconnect losses.
Cross-study evaluation of interfacial stabilization failed because
studies reported initial unmasked voltages rather than quantifying
interfacial bonding enthalpies alongside continuous operational
fill-factor retention under bias.

\subsection{Acknowledgements}\label{acknowledgements}

The authors would like to acknowledge the financial support from Xiamen
University Malaysia through the Xiamen University Malaysia Research Fund
(XMUMRF/2022-C10/IENG/0046), and from National Natural Science
Foundation of China (Ref. No.: 52503390).

\subsection{Declaration of Competing
Interest}\label{declaration-of-competing-interest}

The authors declare no conflict of interest.

\subsection{Data Availability
Statement}\label{data-availability-statement}

The corpus table for the month, the per-paper extraction records with
their verbatim source quotations, the statistics file underlying every
number in the text, and the gate report are provided as Supplementary
Information. Abstract text and full texts are not redistributed.

\subsection{AI Usage Declaration}\label{ai-usage-declaration}

This manuscript was prepared as part of our Research Project for making
a reliable Autonomous Researcher Agent. The agent executed the monthly
review pipeline without human intervention: it harvested the indexed
literature for the month through the OpenAlex, Semantic Scholar,
Crossref and arXiv interfaces, applied a deterministic topical filter,
ranked and selected the papers read closely, extracted every
quantitative claim into a structured record binding each number to a
verbatim quotation from that paper's own abstract, computed the
statistics and figures from those records, and drafted this text from
the resulting evidence briefs. Literature harvesting, selection,
statistics, figures and gate checking are deterministic scripts rather
than model output. Structured extraction used qwen3.8-flash through the
opencode interface; the prose was written by gemini-3.8-flash-high
through the agy interface; independent review used Opus 5 at high
reasoning effort. No figure is a model-generated image: every panel is
plotted directly from the extracted records. Every quantitative
statement in the main text resolves to the machine-generated statistics
file shipped as Supplementary Information, and every citation resolves
to a record carrying its source quotation. The author reviewed and
edited the content, verified the selection, and takes full
responsibility for the content of the published article.

\subsection{References}\label{references}

\begin{enumerate}
\def\labelenumi{\arabic{enumi}.}
\tightlist
\item
  Xu, X. et al.~Noncoplanar multidentate contacts reinforce buried
  interface for perovskite/silicon tandem solar cells. \emph{Science
  Advances} (2026). \url{https://doi.org/10.1126/sciadv.aei2945}
\item
  Wang, F. et al.~Beyond Self‐Assembled Molecules:
  Hydrogen‐Bonding‐Regulated Crystallization of Wide‐Bandgap Perovskite
  for High‐Performance Perovskite‐Silicon Tandem Photovoltaics.
  \emph{Angewandte Chemie International Edition} (2026).
  \url{https://doi.org/10.1002/anie.1676645}
\item
  Liu, Z. et al.~Homogenizing Sn--Pb Distribution Through A‐Site MA
  Cations for Efficient All‐Perovskite Tandem Solar Cells.
  \emph{Advanced Materials} (2026).
  \url{https://doi.org/10.1002/adma.74078}
\item
  Lp, H. et al.~Enhancing Sn-Pb perovskite homogeneity via thioether
  coordination for efficient and stable all-perovskite tandem solar
  cells. \emph{Science Advances} (2026).
  \url{https://doi.org/10.1126/sciadv.aeb8790}
\item
  Xiong, Y. et al.~Radiation-resilient monolithic wide-bandgap
  perovskite/p-type heterojunction silicon tandem solar cells for space
  photovoltaics. \emph{Science Advances} (2026).
  \url{https://doi.org/10.1126/sciadv.aef6600}
\item
  Li, B.E. et al.~Unlocking 27.3\% Perovskite Solar Cells and
  Ultra‐Stable Flexible Modules via Multidentate Molecular‐Mediated
  Buried‐Interface Reconstruction. \emph{Advanced Materials} (2026).
  \url{https://doi.org/10.1002/adma.73807}
\item
  Zhang, Z.Y. et al.~Unlocking 27.3\% perovskite photovoltaics by
  interface-locked dual-molecule contact. \emph{Science Advances}
  (2026). \url{https://doi.org/10.1126/sciadv.aeg1456}
\item
  Gao, D. et al.~n-Type polymer layers enable efficient, scalable, and
  thermally stable perovskite solar modules. \emph{Science} (2026).
  \url{https://doi.org/10.1126/science.aed8175}
\item
  Qian, C. et al.~Bidirectional Chemo-Mechanical Interface Stabilization
  in Perovskite Solar Cells. \emph{Journal of the American Chemical
  Society} (2026). \url{https://doi.org/10.1021/jacs.6c06932}
\item
  Yang, R. et al.~Soft supermolecule stabilized buried interface for
  high-performance inverted perovskite solar cells and modules.
  \emph{Nature Communications} (2026).
  \url{https://doi.org/10.1038/s41467-026-74018-8}
\item
  Shang, X. et al.~A Multisite Molecular Salt Enables Surface
  Passivation and Energetics Management for Efficient Inverted
  Perovskite Solar Cells. \emph{Advanced Energy Materials} (2026).
  \url{https://doi.org/10.1002/aenm.71305}
\item
  Fan, R. et al.~Plasma surface engineering for efficient and stable
  perovskite solar cells and modules. \emph{Science} (2026).
  \url{https://doi.org/10.1126/science.aeg1730}
\item
  Wong, C. et al.~Rational Tailoring of Hole‐Selective Self‐Assembly
  Monolayers Based on Sulfur‐Containing Heterocycles for
  High‐Performance Perovskite Solar Cells. \emph{Angewandte Chemie
  International Edition} (2026).
  \url{https://doi.org/10.1002/anie.1864800}
\item
  Zhang, H. et al.~Dual-Site Interface Passivation by Sulfur-Based Small
  Molecules for the Simultaneous Enhancement of Efficiency and Stability
  in Inverted Perovskite Solar Cells. \emph{ACS Applied Materials \&
  Interfaces} (2026). \url{https://doi.org/10.1021/acsami.6c05331}
\item
  Tan, Y. et al.~Synergistic Interfacial Dangling‐Bond Passivation for
  Enhanced Moisture Resistance in Perovskite Photovoltaics.
  \emph{Advanced Materials} (2026).
  \url{https://doi.org/10.1002/adma.73950}
\item
  Saeed, A. et al.~Amphoteric Anion‐Mediated Crystallization and
  Universal Defect Passivation for High‐Efficiency and Stable Perovskite
  Solar Cells. \emph{Advanced Functional Materials} (2026).
  \url{https://doi.org/10.1002/adfm.77378}
\item
  Chen, R. et al.~Reconstructing the perovskite surface with a
  homogeneous electric potential by reducing localized charge for
  efficient and durable perovskite solar cells. \emph{Energy Materials}
  (2026). \url{https://doi.org/10.20517/energymater.2026.102}
\item
  Jiang, X. et al.~Three-dimensional distribution of amino groups
  improves perovskite crystallization and defect passivation in
  high-performance photovoltaics. \emph{Nature Communications} (2026).
  \url{https://doi.org/10.1038/s41467-026-75990-x}
\item
  Zheng, W. et al.~Intermediate‐Site Anchoring Ligands Enable Robust
  Nonlayered Interfacial Passivation for Efficient and Stable
  Air‐Processed Perovskite Solar Cells. \emph{Advanced Materials}
  (2026). \url{https://doi.org/10.1002/adma.74243}
\item
  Wu, J. et al.~Aniline-Fluoroalkyl Interfacial Passivation Enables
  Defect-Suppressed and Moisture-Resistant Perovskite Solar Cells.
  \emph{ACS Applied Energy Materials} (2026).
  \url{https://doi.org/10.1021/acsaem.6c01560}
\item
  Deng, Z. et al.~Precise Iodide‐Trap Strategy for Photothermally Stable
  Spiro‐OMeTAD‐Based Perovskite Solar Cells. \emph{Angewandte Chemie
  International Edition} (2026).
  \url{https://doi.org/10.1002/anie.2924886}
\item
  Du, X. et al.~Enhancing Interfacial Stability of FAPbI 3 Perovskite
  Solar Cells by Introducing 2D/3D Heterojunction Based on Functional
  Imidazolium Cations. \emph{Advanced Functional Materials} (2026).
  \url{https://doi.org/10.1002/adfm.77015}
\item
  Huang, C. et al.~Phase Separation Suppression by Guanidine Phosphate
  Enables 21.54\% Wide‐Bandgap Perovskite Cells for All‐Perovskite
  Tandem Solar Cells. \emph{Carbon Neutralization} (2026).
  \url{https://doi.org/10.1002/cnl2.70189}
\item
  Cheng, J. et al.~Modulating Growth Kinetics and Defect Passivation in
  Ambient‐Air Evaporation‐Spray Deposition Enables Efficient and Stable
  Rigid/Flexible Perovskite Solar Cells. \emph{EcoEnergy} (2026).
  \url{https://doi.org/10.1002/ece2.70102}
\item
  Lv, D. et al.~Dual‐Site‐Binding Ligand‐Driven (100) Orientation for
  Efficient and Stable Perovskite Solar Cells. \emph{Advanced Materials}
  (2026). \url{https://doi.org/10.1002/adma.74033}
\item
  Fu, J. et al.~Multisite Polymer Regulates Crystallization for
  Low‐Strain Perovskite Film. \emph{Solar RRL} (2026).
  \url{https://doi.org/10.1002/solr.70435}
\item
  Wang, D. et al.~Multi‐Bonding Molecular Interface Bridges Enable
  Synergistic Defect Passivation, Strain Relaxation, and Accelerated
  Electron Extraction in Perovskite Solar Cells. \emph{Advanced
  Functional Materials} (2026). \url{https://doi.org/10.1002/adfm.77373}
\item
  Yang, J. et al.~Dynamic Strain Regulation Via Photoresponsive
  Fullerenes Enables High‐Performance and UV‐Robust Perovskite Solar
  Cells. \emph{Advanced Materials} (2026).
  \url{https://doi.org/10.1002/adma.74220}
\item
  Li, M. et al.~Improved Mechanical Adhesion Enabled by Bilayer Hole
  Transport Layers in Perovskite Solar Cells. \emph{ACS Materials
  Letters} (2026).
  \url{https://doi.org/10.1021/acsmaterialslett.6c00479}
\item
  Wu, H. et al.~A Multidentate Coordination Complex‐Based Self‐Assembled
  Monolayer Stabilizes Perovskite Interface. \emph{Angewandte Chemie
  International Edition} (2026).
  \url{https://doi.org/10.1002/anie.8080769}
\item
  Duan, Y. et al.~Bisphosphonate Self‐Assembled Molecules Enhance
  Anchoring Strength and Perovskite Passivation Toward High‐Efficiency
  Perovskite/Silicon Tandem Solar Cells. \emph{Small} (2026).
  \url{https://doi.org/10.1002/smll.74999}
\item
  Wang, J. et al.~Suppression of SAM desorption by engineering
  polyoxometalate interface for inverted perovskite solar cells.
  \emph{Polyoxometalates} (2026).
  \url{https://doi.org/10.26599/pom.2026.9140147}
\item
  Qin, Z. et al.~Tetraphosphorylated phthalocyanine--based
  self-assembled monolayer stabilizes perovskite photovoltaics.
  \emph{Science Advances} (2026).
  \url{https://doi.org/10.1126/sciadv.aee4175}
\item
  Chen, X. et al.~Molecular Lubricant Mitigates Self‐Assembled Monolayer
  Aggregation in Perovskite Photovoltaics. \emph{Angewandte Chemie
  International Edition} (2026).
  \url{https://doi.org/10.1002/anie.7090517}
\item
  Zhou, B. et al.~Proton Transfer Complex-Confined Self-Assembled
  Monolayer for Thermally Stable Inverted Perovskite Solar Cells.
  \emph{ACS Energy Letters} (2026).
  \url{https://doi.org/10.1021/acsenergylett.6c01489}
\item
  Wan, X. et al.~Strengthened Charge‐Selective Perovskite/SAM
  Heterocontact for Efficient and Durable Perovskite Solar Cells via
  Interfacial Reconfiguration. \emph{Angewandte Chemie International
  Edition} (2026). \url{https://doi.org/10.1002/anie.1018855}
\item
  Su, H. et al.~Synergistic Interface Modulation of Hole Transport for
  Efficient, Stable Inverted Perovskite Solar Cells. \emph{Advanced
  Functional Materials} (2026). \url{https://doi.org/10.1002/adfm.76878}
\item
  Lu, D. et al.~Co‐Assembled Hybrid Interlayer Engineering for Enhanced
  Upper Interface Stability in Inverted Perovskite Solar Cells.
  \emph{Advanced Materials} (2026).
  \url{https://doi.org/10.1002/adma.73886}
\item
  Chen, X. et al.~Perylenocarbazole‐Based Polycyclic Aromatic
  Self‐Assembled Monolayers with Tailored Electrostatic Potentials for
  High‐Performance Organic and Perovskite Solar Cells. \emph{Angewandte
  Chemie International Edition} (2026).
  \url{https://doi.org/10.1002/anie.4031162}
\item
  Ma, P. et al.~Ternary Cooperative Interface With Electronegative
  Phosphine Bridges Enables Efficient and Stable Inverted Perovskite
  Solar Cells. \emph{Advanced Materials} (2026).
  \url{https://doi.org/10.1002/adma.74416}
\item
  Wang, Y. et al.~PFN‐Br Modified Buried Interface Enhanced Charge
  Extraction in Inverted Wide‐Bandgap Perovskite Solar Cells.
  \emph{Small} (2026). \url{https://doi.org/10.1002/smll.74806}
\item
  Miao, Z. et al.~A unified ligand-dimensional design to halt cation
  migration in perovskite photovoltaics. \emph{Science Advances} (2026).
  \url{https://doi.org/10.1126/sciadv.aed6327}
\item
  LIU, M. et al.~In‐Situ Constructed Cations for 2D/3D Perovskite
  Heterostructure for Stable and Efficient Photovoltaics. \emph{Advanced
  Science} (2026). \url{https://doi.org/10.1002/advs.76508}
\item
  Forte, S. et al.~Multilayer Assembly of Perovskite Heterostructures
  via Colloidal Suspensions. \emph{ChemRxiv} (2026).
  \url{https://doi.org/10.26434/chemrxiv.15006632/v1}
\item
  Moazzezi, P. et al.~Solvent Esterification and Stoichiometric Control
  in Ambient-Grown FAPbI 3 Single-Crystal Solar Cells. \emph{Journal of
  the American Chemical Society} (2026).
  \url{https://doi.org/10.1021/jacs.6c05502}
\item
  Zhang, Q. et al.~Sulfolane-Based Solvent Engineering for
  High-Performance MAPbI3 Perovskite Solar Cells. \emph{ACS Applied
  Energy Materials} (2026). \url{https://doi.org/10.1021/acsaem.6c01764}
\item
  Chang, Q. et al.~Dual‐Function Molecular Modulator Enables
  High‐Density Nucleation and Ordered Crystal Growth for Efficient Tin
  Perovskite Solar Cells. \emph{Advanced Functional Materials} (2026).
  \url{https://doi.org/10.1002/adfm.76731}
\item
  Luo, Y. et al.~Impact of Perovskite Surface Termination Dependent
  Molecular Interaction on Interfacial Energetics and Phase Stability.
  \emph{ACS Nano} (2026). \url{https://doi.org/10.1021/acsnano.6c08922}
\item
  Wang, S. et al.~Simultaneous Enhancement of Efficiency and
  Photostability in Wide-Bandgap Perovskite Solar Cells via Managing
  Residual PbI2 with a Phosphoramide Additive. \emph{ACS Applied Energy
  Materials} (2026). \url{https://doi.org/10.1021/acsaem.6c01370}
\item
  Zhu, W. et al.~Thermodynamically Matched Molten‐Salt‐Assisted Phase
  Engineering of 2D/3D Perovskite Heterostructures for Efficient and
  Stable Solar Cells. \emph{Advanced Functional Materials} (2026).
  \url{https://doi.org/10.1002/adfm.76974}
\item
  Jiang, Y. et al.~Structure Optimization Strategy of PbI 2 Vacuum
  Deposition Process for Efficiency Enhancement of Perovskite Solar
  Cells. \emph{Solar RRL} (2026).
  \url{https://doi.org/10.1002/solr.70430}
\item
  Ryu, J. et al.~Resolving Nanoscale Heterogeneities in Lead Halide
  Perovskites Through Low‐Dose Concurrent 4D‐STEM‐EDX Mapping.
  \emph{Advanced Materials} (2026).
  \url{https://doi.org/10.1002/adma.74382}
\item
  Maziviero, F.V. et al.~Computational and experimental evaluation of
  bismuth modification in CsPbX3 (X = I− or Br−): Impact on structural
  and optical properties. \emph{Solar Energy Materials and Solar Cells}
  (2026). \url{https://doi.org/10.1016/j.solmat.2026.114597}
\item
  Zhou, Y. et al.~Preventing Triple Halide Intermediate Formation for
  Segregation-Resistant Wide Band Gap Perovskites. \emph{ACS Energy
  Letters} (2026). \url{https://doi.org/10.1021/acsenergylett.6c01113}
\item
  Maksudov, T. et al.~Achieving 3.31 V in Triple‐Junction
  Perovskite--Organic Tandem Solar Cells via Lattice Distortion--Induced
  Phase Stability. \emph{Advanced Functional Materials} (2026).
  \url{https://doi.org/10.1002/adfm.76715}
\item
  Tan, Q. et al.~Phenomenological Prediction of Halide Perovskite Phase
  Stability Phenomenological Prediction of Halide Perovskite Phase
  Stability via Dielectric-Response Analysis: Equivalence to
  First-Principles Thermodynamics at Five Orders of Magnitude Reduced
  Computational Cost. \emph{ChemRxiv} (2026).
  \url{https://doi.org/10.26434/chemrxiv.15005541/v1}
\item
  Wang, Y. et al.~Interfacial Hydrogen Bonding for Efficient and Robust
  Flexible Tin Perovskite Solar Cells. \emph{Angewandte Chemie
  International Edition} (2026).
  \url{https://doi.org/10.1002/anie.7827890}
\item
  Liu, P. et al.~Sr/Pb Alloying Stabilizes Wide‐Bandgap Perovskites for
  Durable Fully‐Textured Perovskite/Silicon Tandem Solar Cells.
  \emph{Advanced Energy Materials} (2026).
  \url{https://doi.org/10.1002/aenm.71317}
\item
  Ni, Q. et al.~Protonated Cyclohexylamine as an Effective Passivation
  Layer for High-Performance Perovskite Solar Cells. \emph{ACS Applied
  Energy Materials} (2026). \url{https://doi.org/10.1021/acsaem.6c01637}
\item
  Ahmad, N. et al.~Synergistic Interface and Band Engineering for Sn‐Pb
  Perovskites. \emph{Small} (2026).
  \url{https://doi.org/10.1002/smll.74531}
\item
  Zhou, X. et al.~Homogenizing Cesium Distribution via Rubidium
  Incorporation Enables Pure-Iodide 1.67 eV Bandgap Perovskite Solar
  Cells with Efficiency Exceeding 22\%. \emph{ACS Nano} (2026).
  \url{https://doi.org/10.1021/acsnano.6c04084}
\item
  Huang, J. et al.~Multidentate Fullerene Boosts the Thermal Stability
  of Perovskite Solar Cells. \emph{ACS Energy Letters} (2026).
  \url{https://doi.org/10.1021/acsenergylett.6c01517}
\item
  Chalkias, D.A. et al.~Depth‐Resolved Photo‐Degradation of Perovskite
  Layer in Solar Cells Stressed Under ISOS Light‐Aging Protocol
  Conditions: The Role of Substrate, Electron‐Transport‐Layer and
  Perovskite Deposition Method. \emph{Solar RRL} (2026).
  \url{https://doi.org/10.1002/solr.70421}
\item
  Zhong, S. et al.~Molecular Tailoring of Self‐Assembled Monolayers
  Derived From Asymmetric Benzothienoindole for Efficient and Stable
  Inverted Perovskite Solar Cells and Modules. \emph{Angewandte Chemie
  International Edition} (2026).
  \url{https://doi.org/10.1002/anie.1678852}
\item
  Yang, G. et al.~Central Core Engineering of Aromatic Carbonyl
  Molecules Enables Highly Efficient and Stable Perovskite Solar Cells.
  \emph{Angewandte Chemie International Edition} (2026).
  \url{https://doi.org/10.1002/anie.5216345}
\item
  Dai, K. et al.~Thermally Independent Interfacial Reconstruction for
  Phase‐Pure n = 1 2D Perovskites With Mixed‐Orientation Surface
  Architecture. \emph{Advanced Materials} (2026).
  \url{https://doi.org/10.1002/adma.74247}
\item
  You, S. et al.~Boosting Carrier Transport in Carbon‐Based
  Hole‐Transport‐Layer‐Free Perovskite Solar Cells via 2D/3D Interfacial
  Structural Configuration. \emph{Solar RRL} (2026).
  \url{https://doi.org/10.1002/solr.70447}
\item
  Zhao, J. et al.~Eliminating Residue‐Induced Degradation: A Volatile
  Dopant Strategy for High Performance Perovskite Photovoltaics.
  \emph{Small} (2026). \url{https://doi.org/10.1002/smll.74289}
\item
  Kim, S. et al.~Maximum Stability Point Tracking Stabilizes
  Wide‐Bandgap Mixed‐Halide Perovskite Solar Cells. \emph{Advanced
  Energy Materials} (2026). \url{https://doi.org/10.1002/aenm.71358}
\item
  Kong, W. et al.~Stabilizing 1.93-eV ultrawide-bandgap perovskites for
  efficient triple-junction solar cells. \emph{Nature Communications}
  (2026). \url{https://doi.org/10.1038/s41467-026-75305-0}
\item
  Fasti, R. et al.~Dynamic Stress and Strain Evolution and Engineering
  Strategies to Enhance Operational Durability in Perovskite Solar
  Cells. \emph{ACS Energy Letters} (2026).
  \url{https://doi.org/10.1021/acsenergylett.6c01837}
\item
  Guo, B. et al.~A Composition-Dependent Stress--Diffusion Framework
  Identifies Humidity-Tolerant Processing Windows in Hybrid Perovskite
  Solar Cells. \emph{ChemRxiv} (2026).
  \url{https://doi.org/10.26434/chemrxiv-2025-128tn/v2}
\item
  Liu, L. et al.~Artificial Intelligence-Guided Cosolvent Design for
  High-Performance Perovskite/Silicon Tandem Solar Cells.
  \emph{Nano-Micro Letters} (2026).
  \url{https://doi.org/10.1007/s40820-026-02291-9}
\item
  Yun, Y. et al.~Triple-junction all-perovskite solar cells with
  self-assembling hole contacts in all subcells. \emph{Joule} (2026).
  \url{https://doi.org/10.1016/j.joule.2026.102575}
\item
  Dangudubiyyam, U.K. et al.~Eco‐Friendly Solvent System for Inkjet
  Deposition of Wide Bandgap Perovskite Solar Cells Enabling Tandem
  Integration. \emph{Advanced Science} (2026).
  \url{https://doi.org/10.1002/advs.76463}
\item
  Durant, B.K. et al.~In-Operando Thermal Cycling of All-Perovskite
  Tandem Solar Cells. \emph{ACS Energy Letters} (2026).
  \url{https://doi.org/10.1021/acsenergylett.6c00992}
\item
  Chen, Y. et al.~Molecular‐Extrusion‐Driven Halogen Homogenization for
  Efficient Perovskite‐Silicon Tandem Solar Cells. \emph{Angewandte
  Chemie International Edition} (2026).
  \url{https://doi.org/10.1002/anie.5530609}
\item
  Zhou, Y. et al.~Conductive Polymer Bridges Mediate High‐Performance
  Flexible Perovskite--Organic Tandem Solar Cells. \emph{Advanced
  Materials} (2026). \url{https://doi.org/10.1002/adma.74183}
\item
  Harada, N. et al.~Quantitative Hyperspectral Luminescence Imaging
  Reveals Upscaling Losses in Perovskite Solar Modules. \emph{Solar RRL}
  (2026). \url{https://doi.org/10.1002/solr.70375}
\end{enumerate}

\end{document}
